\scanpage{71}\mspage{1}
% archive sheet number: 70

\diagram{\begin{tikzcd}[ampersand replacement=\&] A\arrow[r,shift left,"f"]\&B\arrow[l,shift left,"g"]\end{tikzcd}}

\(g\) adjoint à \deleted{\unsure{droite}}\unsure{gauche} de \(f\), donc on a

\diagram{\begin{tikzcd}[ampersand replacement=\&,column sep=0pt]
|[draw]| \operatorname{id}_A\xrightarrow{\alpha}gf \&
|[draw]| \underset{\substack{\Vert\\\varphi}}{fg}\xrightarrow{\beta}\operatorname{id}_B
\end{tikzcd}}

On a aussi

\diagram{\begin{tikzcd}[ampersand replacement=\&,row sep=small,column sep=small]
\&\&\lambda=f*\alpha*g\\
\boxed{\varphi\xrightarrow{\lambda}\varphi^2}\&\text{i.e.}\&
\deleted{\unsure{gf}}fg\longrightarrow \deleted{\unsure{g}}f\underbrace{gf}g\\
\&\&\operatorname{id}_A\arrow[u,"{\alpha}"']
\end{tikzcd}}

en plus de \(\boxed{\varphi\xrightarrow{\beta}\operatorname{id}_B}\).

\(\lambda\) et \(\beta\) permettent de définir, pour tout objet \(X\) de \(B\), une structure \(1/2\) \deleted{\unsure{covariant|cosimpliciale}} co-semisimpliciale \(\Phi^*(X)\)

\diagram{\begin{tikzcd}[ampersand replacement=\&,column sep=large,row sep=tiny]
\varphi(X)\arrow[r,shift left,"{\deleted{\unsure{\lambda}}\varphi*\beta}"]\arrow[r,shift right,"{\beta*\varphi}"']\arrow[d,equal]\&
\varphi^2(X)\arrow[r,shift left=1.5]\arrow[r]\arrow[r,shift right=1.5]\arrow[d,equal]\&
\varphi^3(X)\arrow[d,equal]\\
\Phi^0(X)\&\Phi^1(X)\&\Phi^2(X)
\end{tikzcd}}

i.e. on a un foncteur co-\(1/2\)-simplicial

\diagram{\begin{tikzcd}[ampersand replacement=\&,column sep=large,row sep=tiny]
\deleted{\unsure{\Phi}}\varphi\arrow[r,shift left,"{\unsure{\varphi*\beta|\varphi\beta}}"]\arrow[r,shift right,"{\beta*\varphi}"']\arrow[d,equal]\&
\varphi^2\arrow[r,shift left=2,"{\varphi^2*\beta\unsure{*u}}"]\arrow[r,"{\varphi*\beta*\varphi}" description]\arrow[r,shift right=2,"{\beta*\varphi^2}"']\arrow[d,equal]\&
\varphi^3\arrow[d,equal]\&\ldots\\
\Phi^0\&\Phi^1\&\Phi^2\&
\end{tikzcd}}

\uline{NB}\quad Les opérations bords n'utilisent que \(\beta\)

\hspace{3em}\rule{4em}{0.4pt}\quad de dégén.\quad\rule{7em}{0.4pt}\(\lambda\)

Pour \unsure{tt|tout} \(Y\) dans \(A\), on a une coaugmentation

\[
f(\deleted{X}Y)\longrightarrow\Phi^*(f(\deleted{X}Y))
\]

qui est une \uline{équivalence d'homotopie}, i.e.

\[
f\xrightarrow{\approx}\Phi^*\circ f
\]

l'homotopie est donnée par

\[
f\xrightarrow{f*\alpha}{\unsure{\Phi|\varphi}}^0\circ f=\varphi f=fgf
\]

Pour tout \(X\) (sous-entendu: on veut exprimer une structure suppl., \deleted{indiquant que} \inserted{\unsure{exprimant si possible} une donnée d'un isom.} \(X\simeq f(Y)\)), on considère les hom. possibles d'objets co\(1/2\) simpliciaux

\[
X\longrightarrow\Phi^*(X)
\]

\diagram{\begin{tikzcd}[ampersand replacement=\&]X\arrow[r,"u"]\&\varphi(X)\arrow[r,shift left,"{\varphi(u)}"]\arrow[r,shift right,"{\lambda}"']\&\varphi^2(X)\end{tikzcd}}

i.e.

\diagram{\begin{tikzcd}[ampersand replacement=\&,row sep=small]
X\arrow[r,dash]\arrow[d,"u"']\&{}\\
\varphi(X)\arrow[r,shift left]\arrow[r,shift right]\&\varphi^2(X)
\end{tikzcd}}
